\section{Data}
\label{sec:data}

The analysis rests on a quarterly panel of eleven euro area countries:
Austria, Belgium, Finland, France, Germany, Greece, Ireland, Italy, the
Netherlands, Portugal and Spain. Starting in 1999, with the common monetary
policy itself, the sample runs to the end of 2025 and therefore covers the
full life of the euro so far: the early convergence years, the housing boom
before 2008, the financial and sovereign debt crises, a decade of very low
rates, and the sharp tightening and subsequent easing of 2022 to 2025. Four
building blocks make up the panel: house prices, the monetary policy shock,
the macroeconomic controls, and the measure of mortgage market structure.
Table \ref{tab:summary} reports summary statistics for the estimation
sample.

\input{tables/summary_stats}

\subsection{House prices}

House prices come from the OECD Analytical House Prices dataset, which I
access through DBnomics. Each country reports a quarterly real residential
property price index with base year 2015. Levels of such an index mean
little across borders, so the analysis works with cumulative log changes
times one hundred, which read as responses in percent and do not depend on
the choice of base year.

\subsection{Monetary policy shock}

Measuring monetary policy is the harder task. Simply using the policy rate
would confuse cause and effect, because the ECB moves that rate in response
to the same economy that drives house prices. The shock series therefore
comes from the Euro Area Monetary Policy event study Database (EA-MPD) of
\citet{altavilla2019}, which records how asset prices move in a narrow
window around each ECB monetary policy event, spanning both the decision
and the press conference. Within such a window little else happens, so the
jump in the three month overnight index swap (OIS) rate captures the
surprise component of the announcement. Compared with the one month rate,
the three month rate is the standard near term policy indicator and is
considerably less noisy.

Not every surprise is a policy shock, however. Sometimes an announcement
mainly reveals what the central bank knows about the economy; interest
rates and stock prices then move in the same direction, and the surprise is
contaminated. Following the sign restriction of \citet{jarocinski2020}, I
compare the sign of each interest rate surprise with the sign of the EURO
STOXX 50 return in the same window. Opposite movements mark a pure monetary
policy shock, and the surprise is kept. Parallel movements mark a central
bank information shock, and the surprise is set to zero. Summing the pure
event level surprises within each quarter then gives the quarterly series.
Without this correction, the raw surprise produces a counterintuitive
positive house price response to tightening, which the cleaned series
removes.

\subsection{Macroeconomic controls}

Three Eurostat series capture the domestic macroeconomic environment. Real
gross domestic product (chain linked volumes, seasonally adjusted) enters
as the quarter on quarter growth rate, and consumer prices enter as the
year on year change in the Harmonised Index of Consumer Prices. For
unemployment I take the monthly harmonised rate and average it to quarters:
the monthly release reaches back to the late 1990s and covers the whole
sample, whereas the quarterly release starts only around 2009 for most
countries. Every control enters with a one quarter lag, so it is
predetermined with respect to the shock. Finally, a dummy marks the
pandemic quarters from 2020 to the first half of 2021, when quarter on
quarter GDP growth reached roughly twenty percent in absolute value and
would otherwise dominate the estimation.

\subsection{Mortgage market structure}

The moderator at the heart of the paper is the structure of the national
mortgage market, that is, whether mortgages are predominantly fixed or
variable rate. Following \citet{badarinza2018} and related euro area
evidence, each country is classified as fixed rate dominant or adjustable
rate dominant. Germany, Austria, France, Belgium and the Netherlands fall
into the first group; Finland, Ireland, Portugal, Spain, Italy and Greece
into the second. A binary split like this is transparent but coarse, since
the variable rate share in fact varies continuously across countries and
over time; I return to this limitation in the concluding section.
